% Section 6.2: 不足与未来工作

\subsection{本文研究的局限}
\begin{itemize}
    \item \textbf{基座单一}：仅 CodeGeeX4-ALL-9B（$\approx$9B），跨规模（1.5B / 14B）与跨家族（DeepSeek-Coder、StarCoder）泛化未验证；
    \item \textbf{数据域单一}：偏算法题，工程型编码任务（API、脚本、框架）未覆盖；
    \item \textbf{$k$ 灵敏度未系统化}：固定 \texttt{FS\_K=2}，更大 $k$ 下的收益曲线及与 \texttt{MAX\_LENGTH} 的相互作用未展开；
    \item \textbf{指标单一}：仅 pass@$k$，未涉及代码质量与生成效率；
    \item \textbf{多 seed 规模有限}：仅锚点组复跑，部分差分量统计显著性不强。
\end{itemize}

\subsection{方法层面}
\begin{itemize}
    \item \textbf{混合推理策略}：若 $\mathrm{H7} \gtrsim \mathrm{H8}$ 且 $\mathrm{H9} \gtrsim \mathrm{H10}$ 同时成立，按题目难度动态切换 Direct / CoT-ZS / CoT-FS；
    \item \textbf{自适应 $k$}：放开「训练/推理 $k$ 匹配」硬约束，按题目侧信号动态调节，混合训练目标覆盖多档示例数；
    \item \textbf{更高效 CoT 目标}：assistant 段内推理/代码加权，或引入 DPO/SimPO 偏好学习。
\end{itemize}

\subsection{系统层面}
\begin{itemize}
    \item \textbf{推理-执行闭环}：子进程执行器接入采样层，基于单测反馈做 self-debug；
    \item \textbf{更大数据规模}：\texttt{HEAD\_N} 扩到 $10^5$ 级并替换更强教师模型，考察 LoRA 增益曲线；
    \item \textbf{工程化部署}：训练-评测流水线封装为在线服务，对接判题与 IDE 插件。
\end{itemize}

\subsection{实验扩展}
本文受限于算力与时间预算，仅以 pass@$k$ 这一聚合层指标得出 $H_1$--$H_{10}$ 之间的差分结论。下列两个方向可在不改动训练与推理主流程的前提下，进一步增强结论的解释力与稳健性：
\begin{itemize}
    \item \textbf{细粒度失效模式归因}：在评测脚本中持久化 per-sample 输出与执行器返回信息，按 Logic / Structural / Runtime / Timeout 四类对失败样例进行自动分类，并辅以典型案例研究。\textit{意义}：将聚合层差分（$\Delta$ pass@$k$）下沉到错误结构层，回答「LoRA 究竟修复了哪一类错误」「CoT 在哪一类题目上反而引入新错误」等机理问题，为后续训练目标与提示策略的针对性优化提供直接证据；
    \item \textbf{鲁棒性消融实验}：在锚点配置（如 H1 / H7 / H10）上对学习率、采样温度、随机种子、\texttt{MAX\_LENGTH} 等单变量做扰动复跑，给出主结论在常见超参邻域内的稳定区间。\textit{意义}：将「H 之间的差分显著性」与「单一扰动源带来的方差」量纲对齐，避免边际收益被超参敏感性所掩盖，使「何时该微调、何时该 CoT、何时该注入示例」的判据具备工程层面的可迁移性。
\end{itemize}
